\section{Related Work}
\label{sec:related_work}

\subsection{Architectures for Biomedical Image Segmentation}
U-Net~\cite{ronneberger2015unet} established the dominant encoder-decoder topology for medical image segmentation, using multi-scale skip connections to preserve high-resolution spatial details. Subsequent variants enhanced representation capacity through nested dense skips (UNet++~\cite{zhou2018unetpp}), additive attention gates (Attention U-Net~\cite{oktay2018attention}), and multi-resolution residual blocks (MultiResUNet~\cite{ibtehaz2020multiresunet}). However, purely convolutional models remain constrained by localized receptive fields, hindering the disambiguation of contiguous anatomical structures.

To overcome this, Vision Transformers (ViTs)~\cite{dosovitskiy2020image} were adopted for global context aggregation. TransUNet~\cite{chen2021transunet} pioneered hybrid architectures by coupling a ResNet-50 CNN with a ViT-B/16 bottleneck transformer. Swin-UNet~\cite{cao2022swinunet} explored a purely transformer-based U-shaped structure with shifted local windows, while UNETR~\cite{hatamizadeh2022unetr} and SwinUNETR~\cite{hatamizadeh2022swinunetr} extended transformers to volumetric 3D segmentation. While effective, these models often exceed 100M parameters, imposing prohibitive memory and computational overhead.

\subsection{Dual-Attention Mechanisms and Computational Bottlenecks}
Dual attention was formalized by Fu et al.~\cite{fu2019dual} (DANet) to simultaneously model spatial dependencies via Position Attention Modules (PAM) and cross-channel interdependencies via Channel Attention Modules (CAM). In biomedical imaging, DA-TransUNet~\cite{sun2024datransunet} embedded dual-attention blocks across the bottleneck and skip connections of TransUNet, establishing state-of-the-art results on multi-organ CT and polyp segmentation. Related efforts explored dual attention in hierarchical transformers (DAE-Former~\cite{azad2023daeformer}, TransDAE~\cite{azad2024transdae}) and state-space architectures (DAMamba~\cite{ding2025damamba}).

However, standard dual attention exhibits quadratic complexity: spatial self-attention scales as $\mathcal{O}(N^2 C)$ for $N = H \times W$ tokens, while channel attention scales as $\mathcal{O}(C^2 N)$. In multi-scale architectures, generating full affinity matrices across multiple high-resolution skip stages ($112 \times 112$ and $56 \times 56$) causes exponential memory consumption and out-of-memory failures on standard clinical workstations.

\subsection{Efficient Attention Approximation and Positioning}
To mitigate quadratic complexity, diverse approximation paradigms have emerged in general vision: local window partitioning (Swin Transformer~\cite{liu2021swin}), low-rank key-value projection (Linformer~\cite{wang2020linformer}, Performer~\cite{choromanski2020performer}), channel group factorizations (EMCAD~\cite{rahman2024emcad}), and efficient transformer backbones (LightViT~\cite{huang2022lightvit}). 

While prior works apply isolated approximation techniques, they retain heavy backbones (such as 86M ViT-B encoders) and high-channel decoders. In our preliminary workshop study~\cite{zhong2026approxda}, we analyzed spatial window sensitivity within the heavy ViT-B baseline. In contrast, this journal work delivers a holistic, system-level lightweight architecture: an early-truncated CNN stem, LightViT-T token processing, multi-scale tri-axial ApproxDA blocks, and a slim $C=64$ decoder, proving that $90\%$ parameter reduction is fully achievable with superior segmentation fidelity across clinical benchmarks.
